\section{Related Work}
\label{sec:related}

\paragraph{Deep-research agent benchmarks measure correctness, not creativity.}
GAIA \citep{mialon2023gaia}, BrowseComp \citep{wei2025browsecomp}, DeepResearch Bench
\citep{du2025deepresearchbench}, LiveResearchBench \citep{liveresearchbench2025},
DeepScholar-Bench \citep{deepscholarbench2025}, and ResearcherBench \citep{researcherbench2025}
score factuality, citation faithfulness, retrieval quality, or report quality. None score whether
a report is creative in the New/Useful/Surprising sense. CREST is complementary: it adds the
missing axis and reuses their citation-faithfulness metrics as a \emph{separate} grounding
check (\S\ref{sec:grounding}), never folding grounding into the creativity score.

\paragraph{Creativity metrics and ideation benchmarks exist, but not for research agents.}
Prior work evaluates LLM idea generation \citep{si2024novel,ruan2024liveideabench}, cross-domain
creativity \citep{houcreativityprism2025}, output diversity \citep{zhang2025noveltybench}, and
linguistic novelty \citep{lu2024creativityindex}. We are careful \emph{not} to claim ``no
creativity test exists.'' Our differentiation is precise: CREST is the first to (i) target
\emph{deep-research agents} that browse and cite, (ii) score the \emph{full} New/Useful/Surprising
triad rather than novelty alone, (iii) organize by \emph{task type} to expose task-dependence, and
(iv) anchor ``new'' to per-item dated known-answers so it reflects synthesis relative to what is
already known, not merely fresher retrieval. That last property is what makes live retrieval
tolerable here; it is \emph{not} backed by a frozen shared retrieval index, which does not exist
(\S\ref{sec:crest}).
Crucially, we avoid n-gram or embedding ``novelty'' metrics that reward text merely for being
unlike anything online; such metrics reward fabrication \citep{saakyan2026death,lu2024creativityindex}.

\paragraph{LLM-as-judge reliability.} LLM judges exhibit self-preference
\citep{panickssery2024own}, assorted biases \citep{ye2024justice}, and can inflate leaderboards
\citep{singh2025leaderboard}; ``novelty'' judgments are especially fragile
\citep{sinhahajari2026mirage}. The literature prescribes a standard mitigation stack
\citep{zheng2023mtbench}: a multi-vendor panel, leave-one-family-out scoring, blindness to entrant
identity, randomised presentation order, inter-judge reliability reporting
\citep{hayes2007alpha,koo2016icc}, and human validation against expert raters. We implement only part of it, and say which part. Implemented: a three-vendor panel
with median aggregation, and reliability reported under both one-way and two-way models. \emph{Not}
implemented: the panel is \textbf{not} order-randomized, \textbf{not} leave-one-family-out (we
measured LOFO and withdrew it---over three vendors it makes judge identity perfectly confounded with
entrant family and \emph{widens} the between-entrant spread), and \textbf{not} human-validated, since
\textbf{no} human study has been run. Blindness is implemented but partially defeated: provider
search tools stamp their own URLs into the report text, leaving a vendor fingerprint in $57\%$ of
records (\S\ref{sec:limitations}). Our own contribution to this literature is the observation that
the stack is insufficient without one further step it does not prescribe---estimating per-judge
\emph{severity}, which on our panel turns out to exceed the whole between-entrant signal.

\paragraph{Accepted benchmark papers set the bar we target, and we do not clear it.} Gorilla,
ToolLLM, MT-Bench \citep{zheng2023mtbench}, DeepScholar-Bench \citep{deepscholarbench2025}, and
MLR-Bench \citep{chen2025mlrbench} were accepted by releasing everything, validating the judge
against humans, running many systems head-to-head, and leaving the benchmark unsolved. CREST meets
the first and last of those and runs five systems head-to-head, but it does \textbf{not} validate its
judge against humans. We state that as a gap rather than a template we follow.
